% ======================================================================
% REVISION CHECKLIST — internal use only; remove before submission.
%
% Usage: every item starts as \citem{ID}{text} (open, empty box).
%   - Mark COMPLETE: change \citem -> \cdone  (green check, grayed text)
%   - Mark REJECTED: change \citem -> \crej   (red box, struck through)
% Requires: amssymb, xcolor, ulem — all already loaded by main.tex.
% To include: \input{revision_checklist} (e.g., right after \maketitle).
% ======================================================================

\newcommand{\citem}[2]{\item[$\square$] \textbf{#1.} #2}
\newcommand{\cdone}[2]{\item[\textcolor{green!55!black}{\rlap{\kern0.25ex\raisebox{0.15ex}{\checkmark}}$\square$}] \textbf{#1.} \textcolor{gray}{#2}}
\newcommand{\crej}[2]{\item[\textcolor{red!70!black}{$\boxtimes$}] \textbf{#1.} \sout{#2}}

\clearpage
\begingroup
\color{red}   % render the entire checklist in red; delete this line (and \endgroup at EOF) to revert
\section*{Revision Checklist (internal --- delete before submission)}

\paragraph{Suggested triage order.} B1--B5 (compile-breaking) $\rightarrow$
A1--A5 (canonical-numbers reconciliation) $\rightarrow$ E1--E4 + abstract
softening $\rightarrow$ C1--C4 $\rightarrow$ rest as time allows. Highest-leverage
decision is A2: settle what the interaction result actually is, then let the
abstract, RQ2/RQ3 results, discussion, and figure captions inherit from it.

\subsection*{A. Critical: internally contradictory numbers}

\begin{itemize}[leftmargin=2em]
\cdone{A1}{Four inconsistent versions of the RQ2 coefficients:
  \texttt{tables/rq2.tex} (Haiku $\hat\beta_H=-32.531$, Sonnet
  $\hat\beta_{PH}=+0.083^{**}$); appendix CI table \texttt{app:logreg\_ci}
  (Haiku $-25.116$, Sonnet $+0.110$); meta-analysis paragraph in
  \texttt{results.tex} (Haiku $+0.029^{**}$, Sonnet $+0.052^{***}$,
  Nano $-0.161$); discussion (Sonnet $+0.083^{***}$, significance differs
  from table's $^{**}$). Also \texttt{results.tex} quotes
  $\hat\beta_H=-25.12$ while describing the table that says $-32.531$.
  \emph{Fix:} pick the canonical fit; regenerate \texttt{rq2.tex},
  \texttt{rq2\_combined.tex}, the CI table, \texttt{rq2\_coefs.png}, and
  rewrite both prose passages from it.}
\cdone{A2}{Headline interaction contradicts itself: Results reports DL pooled
  $\hat\beta_{PH}=+0.037$, CI $[0.015,0.060]$, $p=0.001$; Discussion reports
  $+0.061$, CI $[-0.007,+0.130]$, $p=0.08$ for the same quantity; Sophia's
  note says no significant positive interaction. Adjudicate; if the honest
  summary is ``significant only for Claude Sonnet per-model,'' soften the
  abstract claim to match.}
\cdone{A3}{Claude Sonnet on HLE: in or out? Results text says ``five
  remaining cells'' incl.\ Sonnet $+6.11^\dagger$ and describes its
  trajectory (31\% at $T_4$); \texttt{tab:rq3}, \texttt{tab:coverage}, and
  Limitations all say Sonnet $=\varnothing$ (no $T_0$-correct); Methods says
  only \emph{Haiku} fails on HLE; \texttt{tab:rq4\_hle} reports Sonnet rows
  with Maj.\ Correct exactly 0.500/0.000 (suggesting $n\approx2$--$4$).
  Decide the coverage story once and make methods, results, both tables, and
  limitations agree. Also reconcile $T_0$-correct ranges: ``28--64''
  (results) vs ``28--55'' (limitations).}
\cdone{A4}{Which models have HLE reasoning-trace runs? Results reports
  consensus loss for all six targets; \texttt{app:hle\_consensus} covers only
  GPT-5.4 and Gemini (``the two models with completed pressure runs'');
  Limitations says HLE trace runs exist only for those two. The ``five of six
  targets fragment'' claim rests on data Limitations says does not exist.
  Reconcile (if the appendix table is stale, expand it to all six).}
\cdone{A5}{Does GPT-5.4 have a certain group on GPQA? \texttt{tab:rq1} marks
  Haiku and GPT-5.4 ``N/A --- no certain group''; \texttt{tab:entropy\_dist}
  shows Haiku 22 certain (27.8\%) and GPT-5.4 150 certain (75.8\%) on GPQA.
  Also the \texttt{fig:rq2\_coefs} caption says those two GPQA cells were
  omitted because every question is uncertain, while \texttt{tab:rq2}
  includes both and instead drops GPT-5.4 Mini for separation. Find which
  table reflects the actual sample; fix all three.}
\cdone{A6}{Mathematically impossible mean entropy: methods cites Claude Haiku
  MMLU-Pro $\bar H=-3.21$, but with 10 options $|H|\le\ln 10\approx 2.30$.
  Check log base / whether extra categories (unparsed outputs) enter the
  computation; correct the number.}
\cdone{A7}{Wrong numbers in Discussion's miscalibration paragraph: Haiku GPQA
  confidence ``0.805 $\to$ 0.890'' vs table's 0.630 $\to$ 0.886 (0.805 is the
  MMLU $T_0$ value); ``25--40\,pp'' / ``$\pm$1--3\,pp'' vs results'
  ``20--50\,pp'' / ``$\pm$2--5\,pp''. (Paragraph also duplicates Results ---
  see F8.)}
\cdone{A8}{Appendix single-turn ordering claim falsified by its own table:
  the strict order (expert $>$ social proof $>$ disagreement $>$ accusation
  $>$ crowd $>$ certainty) does \emph{not} hold for all six models (Haiku and
  GPT-5.4 Mini have Disagreement $>$ Social proof). Weaken to what is true,
  e.g.\ ``expert authority is the top category for all six models; personal
  certainty and crowd consensus are always in the bottom two.''}
\cdone{A9}{Gemini's inclusion rule is inconsistent: excluded ex ante from
  MMLU-Pro/GPQA (methods, coverage table) yet included in the single-turn
  MMLU-Pro experiment. Add one sentence: the eligibility criterion gates only
  the uncertainty-conditioned analyses; single-turn category effects do not
  require entropy variance.}
\cdone{A10}{\texttt{tab:rq3} caption says ``HLE rows for Claude Haiku are
  $\varnothing$'' but the table has two $\varnothing$ rows (Haiku and
  Sonnet); ``ClaudeHaiku'' typo; GPQA GPT-5.4 Mini slope of exactly
  $+0.00^\dagger$ is classified ``Capitulator'' --- mark as ``---'' or
  ``Capitulator$^\dagger$'' like the HLE rows.}
\end{itemize}

\subsection*{B. Broken \LaTeX, references, bibliography}

\begin{itemize}[leftmargin=2em]
\cdone{B1}{\texttt{tab:rq4} is undefined: the combined table carrying the
  label is fully commented out in \texttt{tables/rq4.tex}; only
  \texttt{tab:rq4\_mmlu/gpqa/hle} exist, but \texttt{tab:rq4} is referenced
  $\sim$7 times (methods $\times$2, results $\times$2, discussion, appendix
  $\times$2). Update refs to the split tables or restore the label.}
\cdone{B2}{\texttt{tab:coverage} never enters the document: its
  \texttt{\textbackslash input} is commented out in \texttt{methods.tex} but
  it is referenced from results and limitations. Re-enable the input
  (appendix is fine) or drop the refs.}
\cdone{B3}{Missing bib entry \texttt{dersimonian1986meta} (results); also a
  bare \texttt{\textbackslash cite} where the paper otherwise uses
  \texttt{\textbackslash citep}.}
\cdone{B4}{Math typos: the $H_c$ display equation in methods has
  \texttt{\textbackslash bar\{H\}\{\textbackslash mathcal\{M\},\textbackslash
  mathcal\{D\}\}} with a missing subscript underscore (twice); signals (3)
  has $\hat{\gamma}\ell$ and $\hat{c}\ell$ with missing subscript underscores
  (should be $\hat{\gamma}_\ell$, $\hat{c}_\ell$; several occurrences); stray
  $a_i^*$ in the calibration appendix.}
\cdone{B5}{\texttt{\textbackslash usepackage\{ulem\}} without
  \texttt{[normalem]} in \texttt{main.tex} redefines
  \texttt{\textbackslash emph} as underlining --- all italics will render
  underlined. Change to
  \texttt{\textbackslash usepackage[normalem]\{ulem\}}.}
\cdone{B6}{Circular cross-references: signals (1)--(4) each say ``Defined
  formally in Section~2'' pointing at the section they are in. Point to
  \texttt{sec:methods:entropy}, \texttt{:pressure}, \texttt{:calibrator},
  \texttt{:hardness} respectively.}
\cdone{B7}{Wrong ref in results (RQ3 paragraph): ``specification of
  Section \texttt{sec:methods:setup}'' should be \texttt{sec:methods:stats}.}
\cdone{B8}{\texttt{app:logreg\_ci}: undefined symbol $\hat\beta_E$ (elsewhere
  $\hat\beta_H$); caption says GPQA CIs ``available on request'' --- include
  them or write ``omitted for space.''}
\cdone{B9}{Stale duplicate files risk wrong-file edits:
  \texttt{sections/tables/}, \texttt{sections/fig/results\_and\_discussion.tex},
  \texttt{sections/results\_and\_discussion.tex} (not input anywhere), unused
  \texttt{figures/worked\_example.tex}, \texttt{arch\_static.tex},
  \texttt{arch\_with\_adversarial.tex}, \texttt{rq2\_interaction\_combined.png},
  plus \texttt{neurips\_2026.tex}/\texttt{checklist.tex} from another venue.
  Delete or quarantine; confirm \texttt{main.tex} (ACL style) is the build
  target; remove the second commented title; verify \texttt{[review]}
  anonymization.}
\cdone{B10}{Redundant appendix tables: \texttt{rq2\_combined} \emph{and}
  both standalone halves (\texttt{rq2}, \texttt{rq2\_hardness}) are all
  included --- same content three times. Keep the combined one.}
\end{itemize}

\subsection*{C. Missing method details (writable now)}

\begin{itemize}[leftmargin=2em]
\cdone{C1}{Dialogue-history construction under re-polling: with
  $K_{\text{turn}}=5$ re-polls per turn within one dialogue, what assistant
  message enters the transcript for the next turn (one canonical sample?
  modal answer? five parallel histories)? Also state whether RQ4 CoT runs
  share dialogues with the letter-only runs or are separate.}
\cdone{C2}{Two definitions of consensus loss: methods defines it over the
  $R=5$ dialogues' majority beliefs; results (HLE) defines it ``across the
  five CoT samples.'' Pick one, or state that HLE uses the CoT-run variant
  and why.}
\cdone{C3}{Question sampling underspecified: ``stratification by baseline
  entropy bin within each model'' implies per-model subsets, but comparisons
  assume a shared pool; ``198 GPQA questions answered correctly at $T_0$ by
  all included models'' conflicts with the per-model correct-at-$T_0$
  restriction and with $n=79$ for the Claude models in
  \texttt{tab:entropy\_dist}. Explain the actual procedure; state that the
  HLE subset is the multiple-choice portion.}
\cdone{C4}{``Experiment 5'' never defined: results invokes Experiment~5 and
  ensemble confidence $\bar\gamma_{q,t}$, neither in Methods/Experiments.
  Add a short definition (what is ensembled, what ``per-run correctness''
  means, which ECE binning).}
\cdone{C5}{Flip semantics: (i) flips counted only to the planted
  $a^\dagger$ --- say how drifts to third answers are treated; (ii) is the
  per-turn ``flip rate'' in RQ3 state-based (currently on $a^\dagger$) or
  first-flip hazard? The ``resisters recover via self-correction'' reading
  only follows from the state-based version; (iii) the correct-at-$T_0$
  filter (``majority $\hat a_0=a^*$ in at least four of five runs'') mixes
  runs/dialogues terminology --- restate precisely.}
\cdone{C6}{Small omissions: tie-breaking for modal answers; handling of
  unparseable outputs; single-turn per-cell $n$; API access dates / model
  snapshot IDs; explicit note that the calibrator (GPT-5.4 Mini) is also a
  target model (self-evaluation confound; one sentence in Limitations).}
\cdone{C7}{Orphan hyperparameters: \texttt{text-embedding-3-small} and
  ``KMeans clusters (CoT) $=3$'' in \texttt{app:hyperparams} belong to an
  analysis not in the paper. Delete, or reviewers will ask where the
  clustering analysis is.}
\cdone{C8}{Eligibility criterion stated two ways: ``$\geq 2$ populated
  quantile bins'' (methods) vs ``both groups non-empty''
  (\texttt{app:entropy\_dist}). Unify; also fix the pointer to ``the binning
  procedure of Section 2.2'' (bins are defined in Datasets/hyperparams).}
\end{itemize}

\subsection*{D. Statistical framing (addressable in prose)}

\begin{itemize}[leftmargin=2em]
\citem{D1}{Non-independence: turn-level flips are pooled within dialogues
  and questions; Wilson CIs, logistic SEs, and the MWU treat them as
  independent; the mixed model has only a model-level intercept. If feasible
  from saved data, refit with cluster-robust (by-question) SEs; otherwise add
  an explicit caveat in Statistical Models and Limitations. The very narrow
  RQ1 CIs will draw fire without this.}
\citem{D2}{Pressure level is perfectly confounded with rhetorical category
  (fixed order), so $\hat\beta_P$ conflates dose accumulation with category
  potency --- and single-turn results show category effects are
  non-monotonic in the escalation order ($T_4$ strongest). State explicitly,
  where $\hat\beta_P$ is interpreted, that the single-turn experiment is what
  licenses the accumulation reading; soften ``escalating'' to ``fixed-order''
  where it matters.}
\citem{D3}{``Twice the rate'' (abstract, conclusion) is the MMLU-Pro
  cross-model average only; qualify.}
\citem{D4}{Meta-analysis on 5 effect sizes: note Cochran's $Q$ with
  $\mathrm{df}=4$ is underpowered, so non-significant $Q$ is weak evidence of
  homogeneity.}
\citem{D5}{Incomplete significance legend in \texttt{fig:rq2\_coefs} caption:
  ``filled blue square $= p<0.01$, gray cross $=$ n.s.'' --- what marks
  $0.01 \le p < 0.05$?}
\end{itemize}

\subsection*{E. Claim--evidence alignment}

\begin{itemize}[leftmargin=2em]
\cdone{E1}{``Reasoning traces expose capitulation at the first pressure
  turn / earlier than final answer trajectories'' (results header,
  conclusion): no analysis measures lead time between a trace-level signal
  and the final-answer flip; evidence shown is a correctness/confidence
  dissociation plus one worked example. Rephrase to what is shown; cite the
  appendix trace as illustrative. Related: conclusion says belief entropy
  spikes ``at the first pressure turn'' but RQ4 tables skip $T_1$ (and
  $T_3$) --- add the $T_1$ column from existing data or say $T_2$; explain
  the turn subset in the caption.}
\cdone{E2}{Calibrator ``validation'' oversold: methods claims ``we validate
  its reliability post hoc: ECE 0.305, Brier 0.324'' --- poor values, and the
  appendix says $\mathcal{C}$ systematically overstates confidence; also
  82{,}552/ECE 0.305 is just the GPT-5.4-Mini row presented as global.
  Reframe: the calibrator is used for belief extraction ($\hat c_\ell$); its
  confidence is known to be inflated, which is exactly why the analyses use
  the belief--confidence divergence rather than $\hat\gamma_\ell$ alone.}
\cdone{E3}{Calibration appendix conflates two failures: the ``accuracy'' in
  \texttt{tab:cal\_overall}/\texttt{tab:cal\_turn} is whether the target's
  trace supports the gold answer, so falling accuracy reflects target
  capitulation, not calibrator drift --- yet the appendix narrates it as
  ``$\mathcal{C}$ does not register the pressure signal.'' Add a sentence
  acknowledging the two error sources cannot be separated here; align with
  Limitations.}
\cdone{E4}{Abstract/intro scope: ``GPT, Claude, and Gemini families across
  capability tiers'' --- Gemini contributes one model on one benchmark (HLE)
  plus single-turn. Qualify; the abstract currently implies a full
  $6\times3$ grid.}
\cdone{E5}{Dangling concepts: Limitations refers to ``the corrigibility gap
  metric \ldots\ defined but not fully operationalized'' --- never defined in
  the current text; appendix trace refers to ``the categorical reasoning
  shift that Section 7 identifies'' --- Section 7 no longer says this.
  Define briefly or rewrite the pointers.}
\cdone{E6}{``Fragmentation is the default response of contemporary models''
  (discussion) --- hedge given A4's coverage dispute and $n=28$--$72$ per HLE
  cell.}
\end{itemize}

\subsection*{F. Editing / clarity}

\begin{itemize}[leftmargin=2em]
\citem{F1}{Entropy sign-convention contradictions: signals (1) says raw
  entropy ``increases as samples spread'' but \S2.2 defines
  $H=\sum p\log p\le 0$ (decreases); ``Three distinct uncertainty signals''
  introduces four paragraphs, then ``The four signals are complementary.''
  Minimal fix: correct the contradictory sentences and state the convention
  prominently once. (If time allows, flipping to standard $H\ge 0$
  throughout removes a recurring reader stumbling block.)}
\citem{F2}{Terminology drift: ``pressure dose'' vs ``pressure level'' vs
  ``turn index''; ``runs'' vs ``dialogues'' vs ``re-polls''; ``consensus
  loss'' vs ``fragmentation.'' Unify.}
\citem{F3}{Grammar/wording pass: ``we run a six-turn adversarial
  dialogues''; abstract ``In our approach, we investigate\ldots'' and
  ``questions which do not'' $\to$ ``that do not''; conclusion ``We study
  sycophancy with black-box framework'' (missing article), ``We understand
  that pressure level amplifies\ldots'' $\to$ ``We find''; intro ``Existing
  benchmarks explore sycophancy through the flip in stance and its
  frequency'' is awkward.}
\citem{F4}{\texttt{fig:ensemble\_ece} caption self-contradicts: ``ECE is
  near zero at $T_0$ (0.07--0.27)'' (0.27 is not near zero $\to$ ``modest'');
  ``every model moves rightward (confidence unchanged)'' --- if confidence is
  unchanged there is no rightward movement. Rewrite: confidence stays put
  while accuracy falls (downward).}
\citem{F5}{ECE paragraph mixes pooled and per-turn numbers: ``Claude Sonnet
  is the best-calibrated model on MMLU-Pro ($\mathrm{ECE}=0.47$)'' directly
  after saying its $T_0$ ECE is 0.07. State which ECEs are pooled over
  turns.}
\citem{F6}{\texttt{fig:rq2\_interaction}: state that $H=-1.5$ is on the raw
  entropy scale and that curves are predictions from the per-model fits.}
\citem{F7}{Methods \S2.1 overview largely duplicates the formal subsections
  --- if space is needed, shrink to $\sim$6 lines plus the ``three
  resolutions'' sentence.}
\citem{F8}{Deduplicate Results $\leftrightarrow$ Discussion: the
  ``Miscalibration under pressure'' paragraphs appear in both, nearly
  verbatim, with conflicting numbers (A7). Keep numbers in Results; make the
  Discussion paragraph purely interpretive (keep the abstention-rule
  implication).}
\citem{F9}{Empty \texttt{acknowledgements.tex} is input in the review build
  --- remove the input for submission.}
\citem{F10}{Consider re-enabling the commented ``Why the distinction matters
  for deployment'' paragraph in the discussion --- strong framing,
  $\sim$12 lines.}
\end{itemize}
\endgroup